% GENERATED FILE. Edit canonical lessons, then run npm run generate:books.
% canonical-source-hash: fnv1a64:605cb2ba5daf0062
% canonical-lessons: RU-C254-devochka, RU-C254-paren, RU-C254-devushka, RU-C254-roditeli, RU-C254-plemyannik

\chapter{Girl, Lad, Young woman, Parents, Nephew}
\label{ch:ru-girl-lad-young-woman-parents-nephew}
\hlchaptermodality{254}

\begin{chapteropening}
\textbf{By the end of this chapter:} \emph{I can say a girl, a lad, a young woman, parents and a nephew.}

Complete the last lesson of chapter 254: I can say a girl, a lad, a young woman, parents and a nephew.
\end{chapteropening}

\section[\texorpdfstring{\ru{девочка} (dévochka)}{девочка (dévochka)}]{\ru{девочка} (dévochka) — a girl}

\label{lesson:RU-C254-devochka}



\begin{quote}
Before the new one: say the Russian for transport, then the Russian for a trolleybus.
\end{quote}

\subsection*{You'll want to know: \ru{девочка}}
\textbf{\ru{девочка}} (\emph{dévochka}) — \textquotedblleft{}a girl\textquotedblright{}.

\begin{grammarlens}[title={What you've built}]
One more word for this chapter.
\end{grammarlens}

\subsection*{Your turn}
\begin{itemize}
\raggedright
  \item \emph{Say it:} \emph{dévochka}
  \item \emph{Say it:} \emph{dévochka}, once more
  \item \emph{Say it:} say \emph{dévochka}
  \item \emph{Recall:} say the Russian for transport, then the Russian for a trolleybus, then say \emph{dévochka} again
\end{itemize}

\begin{tcolorbox}[breakable,colback=teal!4,colframe=teal!35!black,title={Before you move on}]
What does \ru{девочка} mean? (\textquotedblleft{}A girl\textquotedblright{}.) Say it once more.
\end{tcolorbox}
\section[\texorpdfstring{\ru{парень} (páren)}{парень (páren)}]{\ru{парень} (páren) — a lad, a boyfriend}

\label{lesson:RU-C254-paren}



\begin{quote}
Before the new one: say the Russian for a trolleybus, then the Russian for a girl.
\end{quote}

\subsection*{You'll want to know: \ru{парень}}
\textbf{\ru{парень}} (\emph{páren}) — \textquotedblleft{}a lad, a boyfriend\textquotedblright{}.

\begin{grammarlens}[title={What you've built}]
One more word for this chapter.
\end{grammarlens}

\subsection*{Your turn}
\begin{itemize}
\raggedright
  \item \emph{Say it:} \emph{páren}
  \item \emph{Say it:} \emph{páren}, once more
  \item \emph{Say it:} say \emph{páren}
  \item \emph{Recall:} say the Russian for a trolleybus, then the Russian for a girl, then say \emph{páren} again
\end{itemize}

\begin{tcolorbox}[breakable,colback=teal!4,colframe=teal!35!black,title={Before you move on}]
What does \ru{парень} mean? (\textquotedblleft{}A lad, a boyfriend\textquotedblright{}.) Say it once more.
\end{tcolorbox}
\section[\texorpdfstring{\ru{девушка} (dévushka)}{девушка (dévushka)}]{\ru{девушка} (dévushka) — a young woman, a girlfriend}

\label{lesson:RU-C254-devushka}



\begin{quote}
Before the new one: say the Russian for a girl, then the Russian for a lad.
\end{quote}

\subsection*{You'll want to know: \ru{девушка}}
\textbf{\ru{девушка}} (\emph{dévushka}) — \textquotedblleft{}a young woman, a girlfriend\textquotedblright{}.

\begin{grammarlens}[title={What you've built}]
One more word for this chapter.
\end{grammarlens}

\subsection*{Your turn}
\begin{itemize}
\raggedright
  \item \emph{Say it:} \emph{dévushka}
  \item \emph{Say it:} \emph{dévushka}, once more
  \item \emph{Say it:} say \emph{dévushka}
  \item \emph{Recall:} say the Russian for a girl, then the Russian for a lad, then say \emph{dévushka} again
\end{itemize}

\begin{tcolorbox}[breakable,colback=teal!4,colframe=teal!35!black,title={Before you move on}]
What does \ru{девушка} mean? (\textquotedblleft{}A young woman, a girlfriend\textquotedblright{}.) Say it once more.
\end{tcolorbox}
\section[\texorpdfstring{\ru{родители} (rodíteli)}{родители (rodíteli)}]{\ru{родители} (rodíteli) — parents}

\label{lesson:RU-C254-roditeli}



\begin{quote}
Before the new one: say the Russian for a lad, then the Russian for a young woman.
\end{quote}

\subsection*{You'll want to know: \ru{родители}}
\textbf{\ru{родители}} (\emph{rodíteli}) — \textquotedblleft{}parents\textquotedblright{}.

\begin{grammarlens}[title={What you've built}]
One more word for this chapter.
\end{grammarlens}

\subsection*{Your turn}
\begin{itemize}
\raggedright
  \item \emph{Say it:} \emph{rodíteli}
  \item \emph{Say it:} \emph{rodíteli}, once more
  \item \emph{Say it:} say \emph{rodíteli}
  \item \emph{Recall:} say the Russian for a lad, then the Russian for a young woman, then say \emph{rodíteli} again
\end{itemize}

\begin{tcolorbox}[breakable,colback=teal!4,colframe=teal!35!black,title={Before you move on}]
What does \ru{родители} mean? (\textquotedblleft{}Parents\textquotedblright{}.) Say it once more.
\end{tcolorbox}
\section[\texorpdfstring{\ru{племянник} (plemyánnik)}{племянник (plemyánnik)}]{\ru{племянник} (plemyánnik) — a nephew}

\label{lesson:RU-C254-plemyannik}



\begin{quote}
Before the new one: say the Russian for a young woman, then the Russian for parents.
\end{quote}

\subsection*{You'll want to know: \ru{племянник}}
\textbf{\ru{племянник}} (\emph{plemyánnik}) — \textquotedblleft{}a nephew\textquotedblright{}.

\begin{grammarlens}[title={What you've built}]
One more word for this chapter.
\end{grammarlens}

\subsection*{Your turn}
\begin{itemize}
\raggedright
  \item \emph{Say it:} \emph{plemyánnik}
  \item \emph{Say it:} \emph{plemyánnik}, once more
  \item \emph{Say it:} say \emph{plemyánnik}
  \item \emph{Recall:} say the Russian for a young woman, then the Russian for parents, then say \emph{plemyánnik} again
\end{itemize}

\begin{tcolorbox}[breakable,colback=teal!4,colframe=teal!35!black,title={Before you move on}]
What does \ru{племянник} mean? (\textquotedblleft{}A nephew\textquotedblright{}.) Say it once more.
\end{tcolorbox}
